\chapter*{Feature Engineering for Images (Part 1)}
\addcontentsline{toc}{chapter}{Feature Engineering for Images (Part 1)}

\section*{Image Representation}

Before feeding images into traditional ML models (SVMs, Random Forests, Logistic Regression), we must transform raw pixels into structured, meaningful features. This chapter covers the classical feature engineering pipeline for image data.

\begin{definition}[Image as a Tensor]
A digital image $\mathbf{I}$ is a discrete, multidimensional grid of numerical values:
$$ \mathbf{I} \in \mathbb{R}^{H \times W \times C} $$
where $H$ = height (rows), $W$ = width (columns), and $C$ = number of color channels.
\end{definition}

\begin{description}
    \item[Grayscale:] $C = 1$. Each pixel stores a single intensity value, typically $\in [0, 255]$ (8-bit) where 0 is black and 255 is white.
    \item[Color (RGB):] $C = 3$. Each pixel is a vector $(R, G, B)$ of red, green, and blue intensities. With 8 bits per channel, this yields $256^3 = 16{,}777{,}216$ possible colors.
\end{description}

\begin{lstlisting}[style=pythonstyle]
import cv2
import numpy as np
from PIL import Image

# Load image with OpenCV (returns BGR by default)
img_bgr = cv2.imread("photo.jpg")
img_rgb = cv2.cvtColor(img_bgr, cv2.COLOR_BGR2RGB)
print(img_rgb.shape)   # (H, W, 3) e.g. (480, 640, 3)
print(img_rgb.dtype)   # uint8: values in [0, 255]

# Equivalently with PIL
img_pil = Image.open("photo.jpg")
img_arr = np.array(img_pil)  # Already RGB, shape (H, W, 3)
\end{lstlisting}

\begin{remark}
Raw pixel values are highly sensitive to lighting, translation, rotation, and scale. Flattening an image into $\xx \in \mathbb{R}^{H \cdot W \cdot C}$ and feeding it directly to a model yields poor generalization. We must extract \textbf{invariant} and \textbf{robust} features.
\end{remark}

\section*{Image Metadata (EXIF)}

Many image formats embed \textbf{metadata} via the \textbf{Exchangeable Image File Format (EXIF)} standard, recorded by the capture device. These provide features \textit{without} processing pixel data:

\begin{description}
    \item[Geometric:] Height ($H$), width ($W$), aspect ratio, resolution.
    \item[Photometric:] Color space, bit depth (8-bit vs.\ 16-bit), compression format.
    \item[Acquisition:] Camera model, shutter speed, aperture, ISO, focal length, GPS coordinates, timestamp.
\end{description}

\begin{lstlisting}[style=pythonstyle]
from exif import Image as ExifImage

with open("photo.jpg", "rb") as f:
    meta = ExifImage(f)
if meta.has_exif:
    print(f"Camera: {meta.make} {meta.model}")
    print(f"ISO: {meta.photographic_sensitivity}")
    print(f"GPS: {meta.gps_latitude} {meta.gps_latitude_ref}")
\end{lstlisting}

\section*{Color Spaces and Transformations}

RGB channels are heavily correlated with illumination. Transforming the color space is often the first preprocessing step.

\begin{description}
    \item[Grayscale Conversion:] Reduces $C=3$ to $C=1$, removing color variance when color is task-irrelevant. The standard luminosity formula (ITU-R BT.601) weights channels by human perceptual sensitivity:
    $$ Y = 0.299R + 0.587G + 0.114B $$
    \item[HSV (Hue, Saturation, Value):] Decouples color from brightness.
    \begin{description}
        \item[Hue:] Pure color type as an angle $\in [0^\circ, 360^\circ]$.
        \item[Saturation:] Color purity/intensity.
        \item[Value:] Brightness level.
    \end{description}
    Isolating the \textit{Value} channel provides invariance to shadows and lighting while retaining \textit{Hue} for color-based tasks.
\end{description}

\begin{lstlisting}[style=pythonstyle]
import cv2
import numpy as np

img = cv2.imread("photo.jpg")

# Convert BGR -> Grayscale
gray = cv2.cvtColor(img, cv2.COLOR_BGR2GRAY)
print(gray.shape)  # (H, W) -- single channel

# Convert BGR -> HSV
hsv = cv2.cvtColor(img, cv2.COLOR_BGR2HSV)
hue, sat, val = hsv[:,:,0], hsv[:,:,1], hsv[:,:,2]

# Grayscale intensity range
print(f"Intensity range: [{gray.min()}, {gray.max()}]")
\end{lstlisting}

\section*{Aggregate Statistics (Global Features)}

Treat pixel intensities as a distribution and compute its \textbf{moments}---a compact global summary of brightness and contrast.

\begin{description}
    \item[Mean ($\mu$):] Average brightness.
    $$ \mu = \frac{1}{H \cdot W} \sum_{i=1}^{H} \sum_{j=1}^{W} I(i,j) $$

    \item[Variance ($\sigma^2$):] Intensity spread; correlates with global contrast.
    $$ \sigma^2 = \frac{1}{H \cdot W} \sum_{i=1}^{H} \sum_{j=1}^{W} (I(i,j) - \mu)^2 $$

    \item[Skewness:] Asymmetry of the intensity distribution. Positive skew $\Rightarrow$ tail toward brighter intensities.

    \item[Kurtosis:] ``Peakedness'' of the distribution. High kurtosis $\Rightarrow$ many pixels far from the mean (high contrast/texture regions).
\end{description}

\begin{remark}
Aggregate statistics are computationally cheap and invariant to rotation/translation, but discard \textit{all} spatial information. They are typically used as auxiliary features alongside spatial descriptors (SIFT, HOG).
\end{remark}

\section*{Color Histograms (Global Features)}

A \textbf{color histogram} captures the global distribution of pixel intensities, completely discarding spatial layout.

\begin{description}
    \item[Mechanism:] Divide the intensity range into discrete bins. Count how many pixels fall into each bin.
    \item[Normalization:] Divide by total pixel count to obtain a discrete probability distribution, making the feature invariant to image size.
    \item[Strengths:] Invariant to translation and rotation. Robust to scaling and partial occlusion.
    \item[Weakness:] Complete loss of spatial information---a red apple on green grass produces the same histogram as randomly scattered red and green pixels.
\end{description}

\begin{lstlisting}[style=pythonstyle]
import cv2
import numpy as np

img = cv2.imread("photo.jpg")
gray = cv2.cvtColor(img, cv2.COLOR_BGR2GRAY)

# Compute grayscale histogram: 256 bins over [0, 255]
hist = cv2.calcHist([gray], [0], None, [256], [0, 256])
hist_norm = hist / hist.sum()  # Normalize to probability

# Per-channel color histogram
for i, color in enumerate(["Blue", "Green", "Red"]):
    h = cv2.calcHist([img], [i], None, [256], [0, 256])
    print(f"{color} channel: mean bin count = {h.mean():.1f}")
\end{lstlisting}

\section*{Spatial Patching (Local Features)}

To retain spatial context without the rigidity of full-image flattening, divide the image into a grid of \textbf{sub-regions (patches)} and compute features (histograms, mean intensity, texture statistics) per patch independently. This produces a feature vector that captures \textit{where} certain visual properties occur, bridging the gap between purely global and purely local representations.

\section*{Spatial Filtering and Convolution}

To extract structural features---edges, corners, textures---we apply local neighborhood operations by sliding a small weight matrix across the image.

\begin{definition}[Kernel / Filter]
A \textbf{kernel} $\mathbf{K}$ is a small matrix of weights (typically $3 \times 3$ or $5 \times 5$) that defines how a pixel's neighborhood influences its output value.
\end{definition}

\subsection*{Cross-Correlation vs.\ Convolution}

Let $\mathbf{I}$ be the image and $\mathbf{K}$ a kernel of size $(2k+1) \times (2k+1)$.

\begin{description}
    \item[Cross-Correlation ($\otimes$):] The kernel slides directly over the image without flipping:
    $$ G[i, j] = (\mathbf{I} \otimes \mathbf{K})[i, j] = \sum_{u=-k}^{k} \sum_{v=-k}^{k} \mathbf{K}[u, v] \cdot \mathbf{I}[i+u, j+v] $$
    \item[Convolution ($*$):] The kernel is flipped both vertically and horizontally before multiplication. This ensures associativity in signal processing. In ML, the distinction is often irrelevant since kernels are learned:
    $$ G[i, j] = (\mathbf{I} * \mathbf{K})[i, j] = \sum_{u=-k}^{k} \sum_{v=-k}^{k} \mathbf{K}[u, v] \cdot \mathbf{I}[i-u, j-v] $$
\end{description}

\subsection*{Padding and Stride}
\begin{description}
    \item[Padding:] Filtering shrinks spatial dimensions. To preserve them, add artificial border pixels: \textit{zero-padding} (zeros), \textit{replicate} (copy edge values), or \textit{reflect} (mirror).
    \item[Stride:] The step size of the sliding kernel. Stride $> 1$ produces spatial downsampling.
\end{description}

\subsection*{Common Hand-Crafted Kernels}

\begin{table}[h!]
\centering
\renewcommand{\arraystretch}{1.5}
\begin{tabular}{|c|c|p{7.5cm}|}
\hline
\textbf{Filter} & \textbf{$3 \times 3$ Kernel} & \textbf{Effect} \\
\hline
\textbf{Box Blur} &
$\frac{1}{9} \begin{bmatrix} 1 & 1 & 1 \\ 1 & 1 & 1 \\ 1 & 1 & 1 \end{bmatrix}$ &
Uniform neighborhood average. Low-pass filter: removes high-frequency noise but blurs edges. \\
\hline
\textbf{Gaussian Blur} &
$\frac{1}{16} \begin{bmatrix} 1 & 2 & 1 \\ 2 & 4 & 2 \\ 1 & 2 & 1 \end{bmatrix}$ &
Center-weighted smoothing. More natural, rotationally symmetric blur than box filter. \\
\hline
\textbf{Sharpen} &
$\begin{bmatrix} 0 & -1 & 0 \\ -1 & 5 & -1 \\ 0 & -1 & 0 \end{bmatrix}$ &
Amplifies pixel-neighbor differences. High-pass filter: emphasizes edges and fine details. \\
\hline
\end{tabular}
\end{table}

\begin{lstlisting}[style=pythonstyle]
import cv2
import numpy as np

img = cv2.imread("photo.jpg", cv2.IMREAD_GRAYSCALE)

# Define and apply a sharpening kernel
sharpen_k = np.array([[ 0, -1,  0],
                       [-1,  5, -1],
                       [ 0, -1,  0]], dtype=np.float32)
sharpened = cv2.filter2D(img, -1, sharpen_k)

# Gaussian blur (5x5 kernel, sigma auto-computed)
blurred = cv2.GaussianBlur(img, (5, 5), 0)
\end{lstlisting}

\section*{Edge Detection and Gradients}

Edges are regions of rapid intensity change---they define object boundaries. We detect them by computing the \textbf{discrete derivative} of the image.

\subsection*{Image Gradients}
The gradient of $\mathbf{I}(x,y)$ is a vector of first-order partial derivatives:
$$ \nabla \mathbf{I} = \begin{bmatrix} \frac{\partial \mathbf{I}}{\partial x} \\[4pt] \frac{\partial \mathbf{I}}{\partial y} \end{bmatrix} $$
Since images are discrete, derivatives are approximated via \textbf{finite differences}, implementable as convolution kernels.

\subsection*{Sobel Operators}
The Sobel filters compute the gradient while applying slight Gaussian smoothing to reduce noise:
\begin{align*}
\mathbf{K}_x &= \begin{bmatrix} -1 & 0 & 1 \\ -2 & 0 & 2 \\ -1 & 0 & 1 \end{bmatrix} \quad (\text{detects vertical edges---horizontal gradient}) \\[6pt]
\mathbf{K}_y &= \begin{bmatrix} -1 & -2 & -1 \\ 0 & 0 & 0 \\ 1 & 2 & 1 \end{bmatrix} \quad (\text{detects horizontal edges---vertical gradient})
\end{align*}

After computing $G_x = \mathbf{I} * \mathbf{K}_x$ and $G_y = \mathbf{I} * \mathbf{K}_y$, we derive two per-pixel features:
\begin{description}
    \item[Gradient Magnitude:] Edge strength.
    $$ M = \sqrt{G_x^2 + G_y^2} $$
    \item[Gradient Direction:] Edge orientation (perpendicular to the edge).
    $$ \theta = \arctan\!\left(\frac{G_y}{G_x}\right) $$
\end{description}

\begin{lstlisting}[style=pythonstyle]
import cv2
import numpy as np

img = cv2.imread("photo.jpg", cv2.IMREAD_GRAYSCALE)

# Sobel gradients (64-bit float to capture negative values)
G_x = cv2.Sobel(img, cv2.CV_64F, 1, 0, ksize=3)  # dI/dx
G_y = cv2.Sobel(img, cv2.CV_64F, 0, 1, ksize=3)  # dI/dy

# Gradient magnitude and direction
magnitude = np.sqrt(G_x**2 + G_y**2)
direction = np.arctan2(G_y, G_x)  # Radians in [-pi, pi]
\end{lstlisting}

\begin{remark}
Gradient magnitudes and orientations are the foundation for advanced classical descriptors: \textbf{HOG} (Histogram of Oriented Gradients) and \textbf{SIFT} (Scale-Invariant Feature Transform), which represent the state-of-the-art in pre-deep-learning computer vision.
\end{remark}

\subsection*{Canny Edge Detection Pipeline}
Canny is a multi-stage algorithm for producing clean, thin, well-localized edges:
\begin{enumerate}
    \item \textbf{Noise Reduction:} Apply Gaussian blur to suppress noise.
    \item \textbf{Gradient Calculation:} Compute $M$ and $\theta$ via Sobel operators.
    \item \textbf{Non-Maximum Suppression:} Thin edges by keeping only local maxima along the gradient direction $\theta$---suppresses all non-peak pixels.
    \item \textbf{Hysteresis Thresholding:} Apply two thresholds $T_{\text{high}}$ and $T_{\text{low}}$. Pixels with $M > T_{\text{high}}$ are \textit{strong edges} (always kept). Pixels with $T_{\text{low}} < M < T_{\text{high}}$ are \textit{weak edges}, kept only if connected to a strong edge.
\end{enumerate}

\begin{lstlisting}[style=pythonstyle]
import cv2

img = cv2.imread("photo.jpg", cv2.IMREAD_GRAYSCALE)

# Canny edge detection with two thresholds
edges = cv2.Canny(img, threshold1=100, threshold2=200)
# edges is a binary image: 255 = edge, 0 = non-edge

print(f"Edge pixels: {(edges > 0).sum()}")
print(f"Edge density: {(edges > 0).mean():.3f}")
\end{lstlisting}
